\subsection{Initial topologies of vector spaces}

PROPOSITION 7. --- Let \((\mathrm{E}_{\mathfrak{t}})_{\mathfrak{t}\in\mathbf{I}}\) be a family of topological vector spaces on a topological division ring K. Let E be a vector space on K and for each \(\mathfrak{t}\in I\) , let \(f_{t}\) be a linear mapping of E in \(E_{t}\) . Then the coarsest topology on E which makes each function \(f_{t}\) continuous, is a topology T compatible with the vector space structure of E. Further, if for every \(x\in E\) , \(\phi(x)\) denotes the point \((f_{\mathfrak{t}}(x))\) of the product space \(F=\prod_{\mathfrak{t}\in I}E_{\mathfrak{t}}\) , then the topology T is the inverse image of the topology of the subspace \(\phi(\mathrm{E})\) of F under the linear mapping \(\phi\) .

The last part of the proposition is a particular case of GT, I, § 4.1, prop. 3. The proposition then follows from the next lemma.

Lemma. --- Let M and N be two vector spaces, and g a linear mapping of M in N. If \(T_{0}\) is a topology compatible with the vector space structure of N, then the inverse image of \(T_{0}\) by g is compatible with the vector space structure of M.

We show, for example, that \((\lambda, x) \mapsto \lambda x\) is continuous at each point \((\lambda_{0}, x_{0})\) of \(K \times M\) . Put \(y_{0} = g(x_{0})\) . Every neighbourhood of 0 in M contains a neighbourhood of the form \(g(U)\) where U is a neighbourhood of 0 in N; by hypothesis there exists a neighbourhood V of 0 in K and a neighbourhood W of 0 in N such that the relations \(\lambda - \lambda_{0} \in V\) , and \(y - y_{0} \in W\) imply \(\lambda y - \lambda_{0}y_{0} \in U\) . Thus the relations \(\lambda - \lambda_{0} \in V\) , \(x - x_{0} \in g(W)\) imply \(\lambda x - \lambda_{0}x_{0} \in g(U)\) . We can show similarly that \((x, y) \mapsto x - y\) is continuous in \(M \times M\) .

For each index \(\iota \in I\) , let \(\mathfrak{B}_{\iota}\) be a fundamental system of neighbourhoods of 0 in \(E_{\iota}\) . From the definition of the topology \(\mathcal{T}\) , the filter of neighbourhoods of 0 for this topology is generated by unions of sets of the families \(\overline{f}_{\iota}^{1}(\mathfrak{B}_{\iota})\) ; in other words, the sets of the form \(\bigcap_{k} \overline{f}_{\iota_k}^{-1}(V_{\iota_k})\) form a fundamental system of neighbourhoods of 0 for \(\mathcal{T}\) , the \((\iota_k)_{1 \leqslant k \leqslant n}\) being any finite sequence of indices of I, and, for each index \(k\) , \(V_{\iota_k}\) any set of \(\mathfrak{B}_{\iota_k}\) .

COROLLARY 1. --- Let G be a topological vector space on K. In order that a set H of mappings of G in E be equicontinuous, it is necessary and sufficient that, for all \(\iota \in I\) , the set \(f_{\iota} \circ u\) where u varies in H should be equicontinuous.

This is a particular case of GT, X, § 2.2, prop. 3.

COROLLARY 2. --- If the spaces \(E_{t}\) are Hausdorff, then in order that \(\mathcal{T}\) be Hausdorff, it is necessary and sufficient that, for every \(x \neq 0\) in \(E\) , there should exist an index \(t \in I\) , such that \(f_{t}(x) \neq 0\) .

For \(\phi(E)\) is then a Hausdorff space, and in order that \(\mathcal{T}\) be Hausdorff, it is evidently necessary and sufficient that \(\phi\) be injective; note that we can then identify \(E\) (with \(\mathcal{T}\) ) with the subspace \(\phi(E)\) of \(\prod_{i} E_{i}\) by the mapping \(\phi\) .

COROLLARY 3. --- Suppose the \(E_{t}\) are complete and \(\phi(E)\) is closed in \(F = \prod_{t \in I} E_{t}\) . Then E is complete in the topology T.

For the subspace \(\phi(E)\) of \(F\) is then complete (GT, II, § 3.4, prop. 8 and § 3.5, prop. 10), therefore the same is true of \(E\) in the inverse image topology (GT, I, § 7.6, prop. 10 and GT, II, § 3.1, prop. 4).

* Example. --- Let \(\mathcal{D}'(\mathbf{R})\) be the space of distributions on \(\mathbf{R}\) ; for \(p\) a number such that \(1 \leqslant p \leqslant +\infty\) , let \(j: \mathrm{L}^p(\mathbf{R}) \to \mathcal{D}'(\mathbf{R})\) be the canonical injection, which is continuous (when \(\mathrm{L}^p(\mathbf{R})\) carries its normed space topology and \(\mathcal{D}'(\mathbf{R})\) the strong topology). For every distribution \(f \in \mathcal{D}'(\mathbf{R})\) , denote the derivative of \(f\) by \(\mathrm{D}(f)\) ; recall that \(f \mapsto \mathrm{D}(f)\) is a continuous endomorphism of \(\mathcal{D}'(\mathbf{R})\) . Then let E be the vector subspace of \(\mathrm{L}^p(\mathbf{R})\) formed from those \(f \in \mathrm{L}^p(\mathbf{R})\) for which \(\mathrm{D}(f) \in \mathrm{L}^p(\mathbf{R})\) , and confer on E the coarsest topology making the canonical injection \(i: \mathrm{E} \to \mathrm{L}^p(\mathbf{R})\) and the mapping \(\mathrm{D}: \mathrm{E} \to \mathrm{L}^p(\mathbf{R})\) continuous ( \(\mathrm{L}^p(\mathbf{R})\) carries its normed space topology). For this topology, the space E is complete. For, the image of E in \(\mathrm{F} = \mathrm{L}^p(\mathbf{R}) \times \mathrm{L}^p(\mathbf{R})\) by the mapping \(\phi: f \mapsto (f, \mathrm{D}(f))\) is closed, since it is the trace on \(\mathrm{L}^p(\mathbf{R}) \times \mathrm{L}^p(\mathbf{R})\) of the image G of \(\mathcal{D}'(\mathbf{R})\) in \(\mathcal{D}'(\mathbf{R}) \times \mathcal{D}'(\mathbf{R})\) by the mapping

\[
\phi_ {0}: f \mapsto (f, \mathrm{D} (f));
\]

now \(\mathbf{G}\) is the graph of \(\phi_0\) , therefore closed in \(\mathcal{D}'(\mathbf{R}) \times \mathcal{D}'(\mathbf{R})\) (GT, I, § 8.1, cor. 2 of prop. 2), and as \(\phi(\mathbf{E})\) is the inverse image of \(\mathbf{G}\) by \(i \times i\) , which is continuous, we see that \(\phi(\mathbf{E})\) is closed in F. *

COROLLARY 4. --- Let E be a vector space over a topological division ring K, and let \((\mathcal{T}_{\iota})_{\iota\in I}\) be a family of topologies compatible with the vector space structure of E; then the upper bound T of the topologies \(T_{\iota}\) is compatible with the vector space structure of E.

For, if \(E_{t}\) denotes the topological vector space obtained from E by the topology \(T_{t}\) , and \(f_{t}\) the identity map of E on \(E_{t}\) , then T is the coarsest topology making the \(f_{t}\) continuous.

